% Generated draft: requires mathematical and source review.
\documentclass[11pt]{article}
\usepackage[margin=25mm]{geometry}
\usepackage{fontspec}
\setmainfont{Latin Modern Roman}
\usepackage{amsmath,amssymb,mathtools,unicode-math}
\setmathfont{Latin Modern Math}
\usepackage{enumitem,needspace,xurl,hyperref}
\setlength{\parindent}{0pt}
\setlength{\parskip}{5pt}
\setlength{\emergencystretch}{3em}
\newcommand{\N}{\mathbb{N}}
\newcommand{\Z}{\mathbb{Z}}
\newcommand{\Q}{\mathbb{Q}}
\newcommand{\R}{\mathbb{R}}
\newcommand{\C}{\mathbb{C}}
\newcommand{\F}{\mathbb{F}}
\newcommand{\A}{\mathbb{A}}
\newcommand{\PP}{\mathbb{P}}
\newcommand{\E}{\mathbb{E}}
\newcommand{\eps}{\varepsilon}
\newcommand{\dd}{\,\mathrm{d}}
\begin{document}
\section*{UnsolvedMath ID: 30001164}
\textbf{Source problem code:} catalogue entry
\subsection{Definitions and mathematical statement}
The supplied record bundles several questions and does not specify enough hypotheses to make all of them fully quantified conjectures.

\begin{enumerate}
\item Let \(n\) be a positive integer and let \(V\geq 0\) be a potential on \(\R^n\). Write formally
\[
H_V=-\Delta-V\quad\text{on }L^2(\R^n).
\]
Whenever a self-adjoint realization and a nonempty discrete negative spectrum are justified, denote the negative eigenvalues, counted with multiplicity, by \(E_j\), with \(E_1\) the lowest eigenvalue. The supplied ordering is \(E_1<E_2\leq\cdots\leq0\); only strictly negative eigenvalues enter the sum. For the exponent \(\gamma\), define
\[
S_\gamma(V)=\sum_{E_j<0}|E_j|^\gamma.
\]
The question is whether suitable integrability assumptions on a power of \(V\), strong enough to ensure \(S_\gamma(V)<\infty\), yield
\[
S_\gamma(V)\leq C\,|E_1|^\gamma,
\]
where \(C\) depends only on \(n\), \(\gamma\), and the specified integrability data, uniformly over potentials satisfying those data. The supplied record does not specify the admissible range of \(\gamma\), the power of \(V\), its norm or bound, or the operator hypotheses. Thus the integrability data and the class over which uniformity is requested remain undefined; finiteness of the sum alone does not supply those missing hypotheses.

\item In dimension \(n=1\) and for \(\gamma=1\), the record asks for a direct proof of the Eden--Foias result, or a sharpening by correction terms. It also asks for a direct estimate of the eigenvalue sum without the Birman--Schwinger transformation. No formula for the Eden--Foias inequality, its hypotheses, the proposed correction terms, or the intended direct estimate is supplied. These are therefore retained as research requests whose precise mathematical targets require source clarification.

\item For the geometric question, let \(\mathcal C\) denote the intended class of smooth closed curves of length \(2\pi\) in \(\R^3\). Under the usual regular-curve, arclength interpretation, an admissible curve has a smooth parametrization
\[
r:\R/(2\pi\Z)\longrightarrow\R^3,
\qquad |r'(s)|=1,
\]
and curvature \(\kappa(s)=|r''(s)|\). With periodic boundary conditions, define
\[
L_r=-\frac{\dd^2}{\dd s^2}+\kappa(s)^2,
\qquad
\lambda_1(r)=\inf_{\substack{u\in H^1_{\mathrm{per}}(0,2\pi)\\u\ne0}}
\frac{\displaystyle\int_0^{2\pi}\bigl(|u'(s)|^2+\kappa(s)^2|u(s)|^2\bigr)\,\dd s}
{\displaystyle\int_0^{2\pi}|u(s)|^2\,\dd s}.
\]
The supplied Ovals question asks whether the circle minimizes \(\lambda_1\) in \(\mathcal C\). Since the length-\(2\pi\) circle has curvature \(1\), this is the proposed assertion
\[
\lambda_1(r)\geq1\qquad\text{for every }r\in\mathcal C.
\]
No uniqueness of a minimizer is asserted. The record does not specify whether \(\mathcal C\) requires embeddedness or permits self-intersections or multiple coverings; that choice remains unresolved. Arclength parametrization and periodic boundary conditions make explicit the usual interpretation of the closed-curve operator and require confirmation against the source.
\end{enumerate}

The registry labels this bundle \texttt{partially\_solved} and its statement \texttt{reconstructed\_unverified}. Its supplied verification note expressly withholds full source verification because the integrability exponent and norm are unspecified. These labels and the accompanying literature assessment are supplied curation, not independently verified mathematical conclusions.
\subsection{Short English statement}
This record asks for a uniform comparison of a Schr\"odinger eigenvalue moment with the ground-state energy, direct or improved one-dimensional Eden--Foias estimates, and circle minimization for a curvature-dependent periodic operator. The spectral hypotheses and some geometric conventions remain unspecified, so this is a statement-only draft requiring human mathematical and source review.
\subsection{Sources}
\begin{itemize}
\item Supplied local UnsolvedMath registry record, permanent ID \(30001164\), exact source problem code \texttt{catalogue entry}, titled \emph{Lieb--Thirring Refinements and the Ovals Conjecture}, in the Oberwolfach Reports Open Problems collection. The record gives proposed year 2009 and no individual proposer.
\item Supplied report citation: \emph{Low Eigenvalues of Laplace and Schr\"odinger Operators} (2009), Oberwolfach Reports, pp.\ 355--428, DOI \texttt{10.4171/owr/2009/06}. \url{https://doi.org/10.4171/owr/2009/06}. The record describes these reports as documenting mathematical workshops and their open-problem discussions; it says the bundled questions were reconstructed from this official report.
\item Supplied literature source: Frank--Gontier--Lewin, \emph{Lieb--Thirring inequalities}. \url{https://doi.org/10.1007/s00220-021-04039-5}.
\item Supplied literature source: Benguria--Loss, \emph{Ovals connection}. \url{https://arxiv.org/abs/math-ph/0402048}.
\item Literature assessment dated exactly \texttt{2026-08-21}: the supplied curation labels the status partially solved and states that the one-bound-state conjecture is proved at endpoint exponents and remains open in the intervening one-dimensional range, while the Ovals Problem remains open beyond special families and constraints. This is supplied dated curation derived from the OWR material, not an independently verified assessment. The record explicitly requests independent verification of time-sensitive status and cited claims.
\end{itemize}
\end{document}
